\VerdierPhysicalPageAnchor{0043}%
\VerdierSemanticAnchor{ch1-1-3-8-composition-cont}
of composition in the graded category $\mathcal C[G]$ associated with
the $G^\circ$-category $(\mathcal C,G^\circ,T,c)$ (1.3.2.5).

\VerdierSemanticAnchor{ch1-1-3-9}
\noindent\textbf{1.3.9.}\quad
Let $(\mathcal C,G,T,c)$ be a $G$-category. The category opposite to
$(\mathcal C,G,T,c)$ is the $G^\circ$-category ($G^\circ$ being the
group opposite to $G$)
$(\mathcal C^\circ,G^\circ,T^\circ,c^\circ)$ defined as follows. The
category $\mathcal C^\circ$ is the category opposite to $\mathcal C$.
The functors $T^\circ(g)$, $g\in G^\circ$, are defined by

\VerdierSemanticAnchor{ch1-1-3-9-eq1}
\begin{equation}
  \begin{aligned}
    T^\circ(g)X &= T(g^{-1})X,
      & X&\in\operatorname{Ob}(\mathcal C^\circ)
             =\operatorname{Ob}(\mathcal C),\\
    T^\circ(g)u &= T(g^{-1})u,
      & u&\in\operatorname{Fl}(\mathcal C^\circ)
             =\operatorname{Fl}(\mathcal C),
  \end{aligned}
  \tag{1.3.9.1}
\end{equation}
and the isomorphisms $c^\circ(g,h)$ in $\mathcal C^\circ$ are defined
by

\VerdierSemanticAnchor{ch1-1-3-9-eq2}
\begin{equation}
  \begin{aligned}
    c^\circ(g,h)&\colon
      T^\circ(h\perp g)\longrightarrow T^\circ(h)T^\circ(g),\\
    c^\circ(g,h)&=c^{-1}(g^{-1},h^{-1}).
  \end{aligned}
  \tag{1.3.9.2}
\end{equation}
(we denote composition in $G^\circ$ by $\perp$).

\VerdierSemanticAnchor{ch1-1-3-10}
\noindent\textbf{Remark 1.3.10.}\quad
Let $\mathcal C[G^\circ]\longrightarrow G^\circ$ be the graded category
of type $G^\circ$ associated with the category
$(\mathcal C,G,T,c)$. This category carries a canonical cleavage. Let
$\mathcal C[G^\circ]^\circ\longrightarrow G$ then be the cofibered
category obtained by passing to opposite categories. This category
carries a canonical cocleavage and, since $G$ is a group, it is also
fibered over $G$. It follows that there is a canonical cleavage on
$\mathcal C[G^\circ]^\circ\longrightarrow G$, and hence a
$G^\circ$-category canonically isomorphic to the category opposite to
$(\mathcal C,G,T,c)$, as is verified immediately.

\VerdierSemanticAnchor{ch1-1-3-11}
\noindent\textbf{1.3.11.}\quad
Let $(\mathcal C',G',T',c')$ be a $G'$-category and let
$(\mathcal C'',G'',T'',c'')$ be a $G''$-category. The category
$\mathcal C'\times\mathcal C''$ is canonically endowed with the
structure of a $G'\times G''$-category in the evident way.

\VerdierSemanticAnchor{ch1-1-3-11-product}
The resulting $G'\times G''$-category will be called the
\emph{product} of the $G'$-category $(\mathcal C',G',T',c')$ and the
$G''$-category $(\mathcal C'',G'',T'',c'')$.

\VerdierSemanticAnchor{ch1-1-3-12}
\noindent\textbf{1.3.12.}\quad
Let $(\mathcal C,G^\circ,T,c)$ be a $G^\circ$-category. In practice we
shall most often use right-hand notation for the operations of
$G^\circ$ on $\mathcal C$. For every object $X$ of $\mathcal C$, put

\VerdierSemanticAnchor{ch1-1-3-12-eq1}
\begin{equation}
  T(g)X=X[g]
  \tag{1.3.12.1}
\end{equation}
Likewise, for every morphism $u$ in $\mathcal C$, put

\VerdierSemanticAnchor{ch1-1-3-12-eq2}
\begin{equation}
  T(g)u=u[g].
  \tag{1.3.12.2}
\end{equation}
The isomorphisms $c(g,h)$ therefore define, for every object $X$, the
isomorphisms

\VerdierSemanticAnchor{ch1-1-3-12-eq3}
\begin{equation}
  c(g,h)(X)\colon X[gh]\longrightarrow X[g][h].
  \tag{1.3.12.3}
\end{equation}

\VerdierSemanticAnchor{ch1-1-3-12-notation}
This notation has in particular the advantage of turning the
operations of $G^\circ$ on $\mathcal C$ into right operations of $G$
on $\mathcal C$.
